\documentclass[11pt]{article}
\usepackage[a4paper,margin=26mm]{geometry}
\usepackage{amsmath,amssymb,amsthm}
\usepackage[hidelinks]{hyperref}
\newtheorem{theorem}{Theorem}
\title{Diagonal concavity does not imply global concavity\\
\large A counterexample to MathDB 338076}
\author{Principia Math In Progress}
\date{28 September 2026}
\begin{document}
\maketitle
\begin{abstract}
We disprove Lachaume's diagonal-concavity conjecture with four variables
and strictly positive integer coefficients. The fourth root is strictly
concave on the entire positive diagonal, but has a strictly positive
second directional derivative at a point of the positive orthant.
The proof consists of two explicit polynomial identities.
\end{abstract}

\section*{Task 1: full resolution}
For $x=(x_1,\ldots,x_n)$, let
\[
 \sigma_k(x)=\sum_{1\le i_1<\cdots<i_k\le n}x_{i_1}\cdots x_{i_k},
 \qquad \sigma_0(x)=1.
\]
Conjecture 1 of \cite{Lachaume}, stated on pages 3 and 8 and repeated
on page 18, asserts that for nonnegative coefficients $a_0,\ldots,a_p$,
concavity of
\[
 \left(\sum_{k=0}^p a_k\sigma_k(x)\right)^{1/p}
\]
on the diagonal $x_1=\cdots=x_n>0$ is equivalent to concavity on
$\Gamma_n=(0,\infty)^n$. Here $1/p$-concavity means concavity of the
$1/p$ power, as in the source's Definition 1.2.

\begin{theorem}
Take $n=p=4$ and
\begin{equation}\label{eq:f}
 F(x)=20+40\sigma_1(x)+21\sigma_2(x)+40\sigma_3(x)+20\sigma_4(x).
\end{equation}
Then $t\mapsto F(t,t,t,t)^{1/4}$ has negative second derivative for
every $t>0$, whereas $F^{1/4}$ is not concave on $(0,\infty)^4$.
Consequently $\mathcal K_4^4\ne\mathcal X_4^4$, and the conjecture is false.
\end{theorem}

\begin{proof}
For any positive twice differentiable function $h$ of one real variable,
\begin{equation}\label{eq:calculus}
 (h^{1/4})''=\frac{4hh''-3(h')^2}{16h^{7/4}}.
\end{equation}
On the positive diagonal, the binomial counts in $\sigma_k$ give
\[
 Q(t):=F(t,t,t,t)=20t^4+160t^3+126t^2+160t+20.
\]
Direct differentiation and expansion give the identity
\begin{equation}\label{eq:diagonal}
 3(Q')^2-4QQ''
 =48\bigl[1180(t^2-1)^2+80t(t^2+1)+83t^2\bigr].
\end{equation}
Its right-hand side is strictly positive whenever $t>0$.
Since $Q(t)>0$, formula \eqref{eq:calculus} proves strict concavity
on the whole positive diagonal.

Next set
\[
 x_0=(2,\tfrac12,\tfrac12,\tfrac12),\qquad v=(1,1,1,1),
 \qquad H(u)=F(x_0+uv).
\]
This line lies in the positive orthant for $u>-1/2$.
Substitution into \eqref{eq:f} yields
\begin{equation}\label{eq:line}
 H(u)=20u^4+230u^3+621u^2+713u+\frac{1235}{4}.
\end{equation}
In particular,
\[
 H(0)=\frac{1235}{4},\qquad H'(0)=713,\qquad H''(0)=1242,
\]
and hence
\begin{equation}\label{eq:witness}
 4H(0)H''(0)-3H'(0)^2
 =1235\cdot1242-3\cdot713^2=8763>0.
\end{equation}
Formula \eqref{eq:calculus} now gives
\[
 \left.\frac{d^2}{du^2}F(x_0+uv)^{1/4}\right|_{u=0}
 =\frac{8763}{16(1235/4)^{7/4}}>0.
\]
The restriction of a concave function to any line segment in its domain
is concave, and a twice differentiable concave function has nonpositive
second derivative. This contradiction proves the assertion.
\end{proof}

\paragraph{Scope.}
Every coefficient and every coordinate in the example is strictly
positive. The failure occurs inside the domain, while diagonal concavity
is strict everywhere. No endpoint convention or zero coefficient is used.
One counterexample settles the universally quantified equivalence;
the forward restriction implication remains valid.

The weaker real-rootedness criterion is a distinct question. Indeed,
for $t\ne0$ we have
\[
 \frac{Q(t)}{t^2}=20w^2+160w+86,\qquad w=t+t^{-1}.
\]
The quadratic in $w$ has a root in $(-2,2)$, since its values at $-2$
and $2$ are $-154$ and $486$. The corresponding roots of
$t^2-wt+1$ are nonreal. Thus $Q$ is not real-rooted, and the example
does not challenge the real-rootedness sufficient condition in
\cite[Theorem 2.7(d)]{Gravity}. It is also distinct from the derivative
convolution result previously recorded under MathDB 338079.

\section*{Proposed follow-up tasks}
The original conjecture has been fully refuted above. The following
are four further research tasks, not additional solution claims.
\begin{enumerate}
\setcounter{enumi}{1}
\item Determine whether diagonal $1/3$-concavity implies global
$1/3$-concavity for every $n\ge3$ and every strictly positive
coefficient vector $(a_0,a_1,a_2,a_3)$.
\item Give necessary and sufficient conditions on positive
$(a_0,a_1,a_2,a_3,a_4)$ for
$(\sum_{k=0}^4a_k\sigma_k)^{1/4}$ to be concave on $(0,\infty)^4$.
\item Determine all $b>0$ for which
$(1+2\sigma_1+b\sigma_2+2\sigma_3+\sigma_4)^{1/4}$
is concave on $(0,\infty)^4$; compare this set with the parameters
giving concavity on the positive diagonal.
\item Decide whether for every pair of integers $n\ge p\ge4$ there
are strictly positive coefficients for which
$(\sum_{k=0}^pa_k\sigma_k)^{1/p}$ has negative second derivative
everywhere on the positive diagonal but is not globally concave.
\end{enumerate}

\paragraph{Verification and attribution.}
The accompanying standard-library Python verifier reconstructs both
restrictions from elementary-symmetric sums and checks
\eqref{eq:diagonal}--\eqref{eq:witness} in exact rational arithmetic.
The argument is a complete analytic counterexample, not an inference
from sampled values. This manuscript received author-agent self-review;
independent peer review and proof-assistant verification are not claimed.
The dated literature audit states the public sources checked and its
access limits.

\begin{thebibliography}{9}
\raggedright
\bibitem{Lachaume}
X. Lachaume, \emph{On the concavity of a sum of elementary symmetric
polynomials}, arXiv:1712.10327v1 (2017),
\url{https://arxiv.org/abs/1712.10327}.
\bibitem{Gravity}
X. Lachaume, \emph{The constraint equations of Lovelock gravity theories:
a new $\sigma_k$-Yamabe problem}, arXiv:1712.04528v2 (2018),
\url{https://arxiv.org/abs/1712.04528}.
\end{thebibliography}
\end{document}
